\documentclass[12pt, twoside]{article}
\input{/Users/chris/github/course-files/docs/ib/preamble}

\fancyhead[LO]{La Scuola d'Italia / Dr.\ Huson / IB Math 12 HL \\* 2 November 2026}
\fancyhead[RO]{\fbox{\begin{minipage}{6cm}\footnotesize Nickname:\\[0.5cm]\end{minipage}}}
\fancyhead[LE]{\fbox{\begin{minipage}{6cm}\footnotesize Nickname:\\[0.5cm]\end{minipage}}}
\fancyfoot[C]{\thepage}

\begin{document}
\subsubsection*{5.4 Classwork: Introduction to Differential Equations (no calculator)}

Answer Problems 1 and 2 in the spaces provided. Answer Problems 3--5 on
lined paper. Show all working and give exact answers.

\begin{enumerate}[itemsep=0.15cm]

%--- Problem 1: Verify a general solution; particular solution; equilibrium value; form a DE (Newton's law of cooling) [9 marks] ---
%--- Source: Original (freshly authored 2026-11-02, HL syllabus AHL 5.18) ---
\item[1.] [Maximum mark: 9]

Consider the differential equation
$$\frac{\mathrm{d}y}{\mathrm{d}x} = 6 - 2y.$$

\begin{enumerate}[itemsep=0.15cm]
  \item Show that $y = Ae^{-2x} + 3$ is a solution of the differential
  equation for every constant $A$. \hfill [3]
  \item Find the particular solution for which $y = 5$ when $x = 0$.
  \hfill [2]
  \item Write down the value that $y$ approaches as $x \to \infty$, and
  show that this value of $y$ makes $\dfrac{\mathrm{d}y}{\mathrm{d}x} = 0$.
  \hfill [2]
\end{enumerate}

A cup of tea cools in a room at $20^{\circ}\mathrm{C}$. The rate at which
its temperature $T$ decreases is proportional to the difference between
$T$ and the room temperature.

\begin{enumerate}[itemsep=0.15cm]
  \item[(d)] Write down a differential equation for $T$ in terms of the time
  $t$ and a positive constant $k$. \hfill [2]
\end{enumerate}

\begin{tikzpicture}
  \draw (-0.6,0) rectangle (14.6,10.0);
  \foreach \y in {1,2,...,9} \draw[dotted] (0,\y)--(14.1,\y);
\end{tikzpicture}

\newpage

%--- Problem 2: Separation of variables: dy/dx = 6x^2 y with y(0) = 2; dy/dx = e^(x - y) with y(0) = ln 2 [9 marks] ---
%--- Source: Original (freshly authored 2026-11-02, HL syllabus AHL 5.18) ---
\item[2.] [Maximum mark: 9]

Solve each differential equation by separating the variables. Give each
answer in the form $y = f(x)$.

\begin{enumerate}[itemsep=0.15cm]
  \item $\dfrac{\mathrm{d}y}{\mathrm{d}x} = 6x^{2}y$, \ given that $y = 2$
  when $x = 0$. \hfill [5]
  \item $\dfrac{\mathrm{d}y}{\mathrm{d}x} = e^{x - y}$, \ given that
  $y = \ln 2$ when $x = 0$. \hfill [4]
\end{enumerate}

\begin{tikzpicture}
  \draw (-0.6,0) rectangle (14.6,17.4);
  \foreach \y in {1,2,...,16} \draw[dotted] (0,\y)--(14.1,\y);
\end{tikzpicture}

\newpage

\emph{Answer Problems 3--5 on lined paper.}

%--- Problem 3: Newton's law of cooling: solve dT/dt = -k(T - 20), T(0) = 90; k from T(10) = 55; time to reach 37.5 [9 marks] ---
%--- Source: Original (freshly authored 2026-11-02, HL syllabus AHL 5.18); continues Problem 1(d) ---
\item[3.] [Maximum mark: 9]

The tea in Problem 1 is $90^{\circ}\mathrm{C}$ when $t = 0$, where $t$ is
the time in minutes. Its temperature satisfies
$$\frac{\mathrm{d}T}{\mathrm{d}t} = -k\left(T - 20\right), \qquad k > 0.$$

\begin{enumerate}[itemsep=0.15cm]
  \item By solving the differential equation, show that
  $T = 20 + 70e^{-kt}$. \hfill [4]
  \item After $10$ minutes the temperature of the tea is
  $55^{\circ}\mathrm{C}$. Show that $k = \dfrac{\ln 2}{10}$. \hfill [2]
  \item Find the time at which the temperature of the tea is
  $37.5^{\circ}\mathrm{C}$. \hfill [3]
\end{enumerate}

\bigskip

%--- Problem 4: dy/dx = x + y, y(0) = 1: Euler h = 0.5 to x = 1; exact solution by integrating factor; concavity explains the underestimate [10 marks] ---
%--- Source: Original (freshly authored 2026-11-02, HL syllabus AHL 5.18) ---
\item[4.] [Maximum mark: 10]

Consider the differential equation
$\dfrac{\mathrm{d}y}{\mathrm{d}x} = x + y$, where $y = 1$ when $x = 0$.

\begin{enumerate}[itemsep=0.15cm]
  \item Use Euler's method with a step length of $0.5$ to find an
  approximate value of $y$ when $x = 1$. \hfill [3]
  \item By using an integrating factor, solve the differential equation.
  Give your answer in the form $y = f(x)$. \hfill [5]
  \item Show that $\dfrac{\mathrm{d}^{2}y}{\mathrm{d}x^{2}} = 1 + x + y$.
  Hence explain why the approximation in part (a) is less than the true
  value of $y$ when $x = 1$. \hfill [2]
\end{enumerate}

\newpage

%--- Problem 5: Homogeneous DE dy/dx = (x^2 + y^2)/(xy), substitution y = vx, y(e) = 2e [8 marks] ---
%--- Source: Original (freshly authored 2026-11-02, HL syllabus AHL 5.18) ---
\item[5.] [Maximum mark: 8]

Consider the differential equation
$$\frac{\mathrm{d}y}{\mathrm{d}x} = \frac{x^{2} + y^{2}}{xy},
\qquad x > 0,\ y > 0.$$

\begin{enumerate}[itemsep=0.15cm]
  \item Use the substitution $y = vx$ to show that
  $x\dfrac{\mathrm{d}v}{\mathrm{d}x} = \dfrac{1}{v}$. \hfill [3]
  \item Given that $y = 2e$ when $x = e$, solve the differential equation.
  Give your answer in the form $y = f(x)$. \hfill [5]
\end{enumerate}

\end{enumerate}
\end{document}
